\section{Дослідження генерації простих чисел}

\subsection{Для чого при генерації ключів RSA необхідні великі прості числа?}
а
 
\subsection{Чому на практиці для великих чисел використовуються ймовірнісні тести простоти?}

\subsection{У чому полягає основна ідея тесту Ферма та чому позитивний результат цього тесту не гарантує простоту числа?}

\subsection{Яку проблему тесту Ферма частково вирішують більш досконалі ймовірнісні тести?}

\subsection{Який принцип роботи тесту Міллера--Рабіна та яке значення має кількість його ітерацій?}

\subsection{Чому дуже мала ненульова ймовірність помилки тесту простоти може бути прийнятною при практичній генерації криптографічних ключів?}

\subsection{Які вимоги, крім достатньої довжини, висуваються до простих чисел, що використовуються при формуванні RSA-ключів?}

!!!коротко порівняти